\section{Pfaffian functions and Pfaffian sets}\label{sec:pfaffian}
References for the material in this section are~\cite{khovanskii1991fewnomials, zell2003quantitative, gabrielov2004complexity}. We assume throughout that $n\in \N_{>0}$. Pfaffian functions are defined relative to an open domain $\cU\subseteq\R^n$; unless a domain is specified we take $\cU=\R^n$, which covers the neural network activation functions of interest.

\subsection{Pfaffian functions}
Pfaffian functions, introduced by Khovanskii~\cite{khovanskii1991fewnomials}, are functions that satisfy triangular systems of first-order partial differential equations with polynomial coefficients. For chains defined on all of $\R^n$, the sets defined by Pfaffian functions are tame in the sense of o-minimal geometry~\cite{wilkie1996model, wilkie1999theorem, speissegger1999pfaffian}, and many of their geometric and topological properties, such as the number of connected components and the sum of Betti numbers, can be bounded effectively in terms of the format~\cite{gabrielov2004complexity}.

\begin{definition}[Pfaffian function]
Let $\cU \subset \R^n$ be an open set. A \emph{Pfaffian chain} of length (also called order) $s \in \N$ and chain degree $\alpha \in \N_{>0}$ over $\cU$ is a sequence of functions $\boldq= (q_1, \dotsc, q_s)$ with $q_i\in C^\infty(\cU)$ for $i\in [s]$, such that there exist polynomials $P_{ij}\in \R[X_1,\dots,X_n,Y_1,\dots,Y_i]_{\leq \alpha}$, for $i\in [s]$ and $j\in [n]$, that verify
\begin{equation}\label{eq:pfaffchain}
\frac{\partial q_i}{\partial x_j}(x) = P_{ij}(x, q_1(x), \dotsc, q_i(x)).
\end{equation}
A function $g(x)=P(x,q_1(x),\dots,q_s(x))$, with $P\in \R[X_1,\dots,X_n,Y_1,\dots,Y_s]_{\leq \beta}$ for $\beta\in \N$, is called a \emph{Pfaffian function} of chain degree $\alpha$, degree $\beta$, and chain length $s$. A function $\cU\to \R^m$ is called Pfaffian if all its components are Pfaffian.
\end{definition}

The triple $(\alpha, \beta, s)$ is called a \emph{format} of $g$; its degree parameters are upper bounds for the chosen representation. We denote by $\Pfaff_{\boldq,\beta}(\cU)$ the set of all Pfaffian functions over $\cU$ with chain $\boldq$ and degree at most $\beta$, and by $\Pfaff_{\alpha,\beta,s}(\cU)$ the set of all Pfaffian functions over $\cU$ with format $(\alpha,\beta,s)$. When a map $\cU\to \R^m$ has a particular format, every component admits that format, and we write $\Pfaff_{\alpha,\beta,s}(\cU; \R^m)$. A common chain for the components is required only when explicitly stated.

\begin{remark}
 Note that the Pfaffian chain associated with a Pfaffian function, and hence also a format, is not unique. In particular, $\Pfaff_{\alpha,\beta,s}(\cU)\subseteq \Pfaff_{\alpha',\beta',s'}(\cU)$ for $\alpha'\geq \alpha$, $\beta'\geq \beta$, and $s'\geq s$. In the following, we will often leave the chain implicit, and our results will be stated in terms of the format.
\end{remark}

We call a Pfaffian function $g\colon \cU\to \R$ \emph{autonomous} with respect to a Pfaffian chain $\boldq$ if 
\begin{equation}\label{eq:simplepfaff}
  \frac{\partial q_i}{\partial x_j}(x) = P_{ij}(q_1(x),\dots,q_i(x))
\end{equation}
for some $P_{ij}\in \R[Y_1,\dots,Y_i]$, and $g(x)=P(q_1(x), \dots, q_s(x))$ for some $P\in \R[Y_1,\dots,Y_s]$. 
Note that being autonomous is a property of a given Pfaffian representation of a function, not a property of the function itself. Every Pfaffian function can be made autonomous by adjoining the coordinate functions $x_1,\dots,x_n$ to the Pfaffian chain. In what follows, we often work with a fixed chain that may be implicit, and refer to a Pfaffian function as autonomous if it is autonomous with respect to that chain.

\begin{example}\label{ex:1} The following are simple examples of Pfaffian functions.
        \begin{enumerate}
            \item A polynomial $P\in\R[X_1,\dots,X_n]_{\le\beta}$ gives a Pfaffian function with format $(\alpha,\beta,0)$ for any $\alpha>0$;
            \item $\exp\colon \R\to \R$ is Pfaffian with format $(1,1,1)$. A Pfaffian chain is $\boldq=(\exp)$ and $P_{11}(X,Y_1)=Y_1$;
            \item $\tanh$ is Pfaffian with format $(2,1,1)$. A Pfaffian chain is $\boldq=(\tanh)$ and $P_{11}(X,Y_1)=1-Y_1^2$;
            \item The logistic sigmoid $\sigma(x) = (1 + \econst^{-x})^{-1}$ is Pfaffian with format $(2, 1, 1)$. A Pfaffian chain is $\boldq=(\sigma)$ and $P_{11}(X,Y_1)=Y_1(1-Y_1)$;
            \item $\arctan$ is Pfaffian with format $(3, 1, 2)$. A Pfaffian chain is $\boldq=((1+x^2)^{-1},\arctan(x))$. The polynomials are $P_{11}(X,Y_1) = -2XY_1^2$ and $P_{21}(X,Y_1,Y_2)=Y_1$. 
        \end{enumerate}
Examples (2--4) are autonomous with respect to the given chains, whereas nonconstant polynomials in (1) and the arctangent representation in (5) are not.
\end{example}

%\begin{definition}[Algebraically independent Pfaffian chain]
%\label{defn:algebraically-independent-pfaffian-chain}
%A Pfaffian chain $\boldq = (q_1, \ldots, q_r)$ is called algebraically independent if for all non-zero $P \in \R[Y_1, \ldots, Y_r]$, $P(q_1(x), \ldots, q_r(x))$ is not identically zero.
%\end{definition}

%\begin{remark}
%All the cases in Example~\ref{ex:1} have algebraically independent Pfaffian chains. In~\cite{lotz2024partitioning} it is shown that the chain $\boldq=(1/x,x^m)$ is algebraically independent if $m\in \R\backslash \mathbb{Q}$.    
%\end{remark}

Pfaffian functions enjoy some convenient closure properties under algebraic operations and composition, as illustrated in the following two results (see also~\cite[Proposition 1.8]{zell2003quantitative}).

\begin{lemma}\label{le:pfaffalgebra}
Let $g\in \Pfaff_{\alpha, \beta, s}(\cU)$ and $h\in \Pfaff_{\alpha', \beta', s'}(\cU)$, with $\beta,\beta'\ge1$. Then:
\begin{enumerate}
\item $g + h \in \Pfaff_{\max\{\alpha,\alpha'\}, \max\{\beta,\beta'\},s+s'}(\cU)$;
\item $gh \in \Pfaff_{\max\{\alpha,\alpha'\}, \beta+\beta',s+s'}(\cU)$;
\item For each $i\in [n]$, $\partial_i g \in \Pfaff_{\alpha,\alpha+\beta-1,s}(\cU)$, with respect to the same chain as $g$.
\end{enumerate}
\end{lemma}

\begin{proof}
The first two statements are straightforward consequences of the fact that the concatenation of Pfaffian chains is again a Pfaffian chain. For the third statement, assume $g(x)=P(x,q_1(x),\dots,q_s(x))$ for a polynomial $P\in \R[X_1,\dots,X_n,Y_1,\dots,Y_s]_{\leq \beta}$. Then
\begin{align*}
\partial_i g(x)={}&(\partial_{X_i}P)(x,\boldq(x))\\
&+\sum_{\ell=1}^s(\partial_{Y_\ell}P)(x,\boldq(x))
P_{\ell i}(x,q_1(x),\dots,q_\ell(x)).
\end{align*}
The resulting expression is a polynomial in $x$ and the $q_j$, with degree bounded by $\alpha+\beta-1$.
\end{proof}

\begin{remark}\label{re:1}
The chain-length bounds for $g+h$ and $gh$ are in general not sharp: if $g$ and $h$ share a chain $\boldq$ of length $s$, then their sum and product use that same chain. Likewise, for a constant vector $\tau$, $\partial_\tau g=\sum_i\tau_i\partial_i g$ has format $(\alpha,\alpha+\beta-1,s)$ on the chain of $g$. More generally, a linear combination of Pfaffian functions $g_1,\dots,g_m$ with formats $(\alpha_i,\beta_i,s_i)$ has format $(\max_i\{\alpha_i\},\max_i\{\beta_i\},\sum_i s_i)$.
\end{remark}

\begin{lemma}\label{le:composition}
Let $g\in \Pfaff_{\alpha,\beta,s}(\cU;\R^m)$ and $h\in \Pfaff_{\alpha',\beta',s'}(\R^m)$ be Pfaffian functions, and assume $\beta,\beta'\ge1$ and $s'\geq 1$.
\begin{enumerate}
\item $h\circ g \in \Pfaff_{(\alpha'+1)\beta+\alpha-1, \beta\beta', ms+s'}(\cU)$;
\item If $h$ is autonomous, then $h\circ g \in \Pfaff_{\alpha'+\alpha+\beta-1,\beta', ms+s'}(\cU)$;
\item If all components of $g$ share a Pfaffian chain $\boldq$ of length $s$, then $h\circ g$ admits a chain of length $s+s'$, with the degree bounds in (1) or (2), respectively.
\end{enumerate}
\end{lemma}

\begin{proof}
Let $g=(g_1,\dots,g_m)$, and for each $i\in [m]$, let $\boldq^{(i)}=(q_1^{(i)},\dots,q_s^{(i)})$ be a Pfaffian chain for $g_i$. Let $\boldq'=(q_1',\dots,q_{s'}')$ be a Pfaffian chain for $h$. Let $P$ and $P_i$, $i\in [m]$, be the polynomials representing $h$ and the $g_i$, respectively. The composition $h\circ g$ is then represented as 
\begin{equation}\label{eq:longpoly}
 (h\circ g)(x) = P(P_1(x,\boldq^{(1)}(x)),\dots,P_m(x,\boldq^{(m)}(x)),\boldq'(g(x))),
\end{equation}
which is clearly a polynomial in $x=(x_1,\dots,x_n)$ and in the chain
\begin{equation}\label{eq:newchain}
  (q_1^{(1)},\dots,q_s^{(1)},\dots,q_1^{(m)},\dots,q_s^{(m)},q'_1\circ g,\dots,q_{s'}'\circ g)
\end{equation}
of length $ms+s'$. It remains to be seen that this is a Pfaffian chain. For the first $ms$ functions in this chain there is nothing to show. It remains to be seen that the derivatives of $q_i'\circ g$ can be expressed as polynomials in $x$ and in the previous elements of the chain. In fact, for $j\in [n]$ we have
\begin{equation}\label{eq:composite}
\begin{aligned}
\partial_j(q'_i\circ g)(x)
&=\sum_{\ell=1}^m(\partial_\ell q'_i)(g(x))\,\partial_j g_\ell(x)\\
&=\sum_{\ell=1}^m P'_{i\ell}(g(x),q'_1(g(x)),\dots,q'_i(g(x)))\,\partial_j g_\ell(x),
\end{aligned}
\end{equation}
where the $P_{i\ell}'$ are the polynomials associated to the chain $\boldq'$. 
This expresses the partial derivative as a polynomial in the coordinates $x$ and the chain elements in~\eqref{eq:newchain} up to $q_i'\circ g$. Indeed, $g_{\ell}$ and, by Lemma~\ref{le:pfaffalgebra}(3), its partial derivatives are polynomials in $x$ and $q_1^{(\ell)},\dots,q_s^{(\ell)}$. This establishes that~\eqref{eq:newchain} is a Pfaffian chain of length $ms+s'$.

To prove the degree bounds in the general case (1), note that the degree of the composition of polynomials is bounded by the product of the degrees. The expression $P_{i\ell}'(g(x),q'_1(g(x)),\dots,q'_i(g(x)))$ has degree at most $\alpha'\beta$ in $x$ and the new chain elements. Lemma~\ref{le:pfaffalgebra}(3) bounds the degree of $\partial_j g_\ell$ by $\alpha+\beta-1$, establishing the bound $(\alpha'+1)\beta+\alpha-1$ on the degree of the new Pfaffian chain~\eqref{eq:newchain}.
Moreover, the degree of the polynomial~\eqref{eq:longpoly} representing $h\circ g$ is clearly bounded by $\beta\beta'$. 
For the autonomous case (2), the polynomials $P'_{i\ell}$ have no explicit dependence on the input coordinates. Thus $P'_{i\ell}(q'_1(g(x)),\dots,q'_i(g(x)))$ has degree at most $\alpha'$ in the last $s'$ chain elements, while $\partial_j g_\ell$ has degree at most $\alpha+\beta-1$ in the remaining variables. This bounds the degree of~\eqref{eq:composite} by $\alpha'+\alpha+\beta-1$. We also observe from~\eqref{eq:longpoly} that in the autonomous case, $h\circ g$ depends only on the last $s'$ elements of its Pfaffian chain in the same way that $h$ depends on its Pfaffian chain $\boldq'$. This implies that the degree is bounded by $\beta'$.
Case (3) is clear.
\end{proof}

\begin{remark}\label{re:affine}
Pullback by an affine map preserves a Pfaffian format: substituting affine expressions in~\eqref{eq:pfaffchain} does not increase the polynomial degrees. In particular, a rigid change of coordinates and restriction to an affine subspace preserve the format. This also covers polynomial functions represented by an empty chain.
\end{remark}

We are also interested in the geometric objects defined by Pfaffian functions.

\begin{definition}
A \emph{Pfaffian set} (or Pfaffian variety) is a set of the form 
\begin{equation*}
  V = \mathcal{Z}(g_1, \dotsc, g_k) = \{x\in \R^n\colon g_1(x)=\cdots = g_k(x)=0\},
\end{equation*}  
where the $g_i\colon \R^n\to \R$ are Pfaffian functions. A {\em basic semi-Pfaffian set} is a set of the form
\begin{equation*}
 B = \{x\in \R^n \colon g_1(x)=\cdots = g_k(x)=0, h_1(x)>0,\dots,h_{\ell}(x)>0\},
\end{equation*}
where the $g_i$ and $h_j$ are Pfaffian functions. A {\em semi-Pfaffian set} is a finite union of basic semi-Pfaffian sets.
\end{definition}

Note that semi-Pfaffian sets are precisely the sets that can be written as unions, intersections, and complements of sets defined by expressions of the form $g_i(x) \star 0$, with $\star \in \{=,\, \leq,\, \geq\}$. 

\begin{remark}
  Since the concatenation of Pfaffian chains is again a Pfaffian chain, we can assume without loss of generality that the functions $g_i$ and $h_j$ in the definition of a (semi-)Pfaffian set are Pfaffian with respect to the same Pfaffian chain.
\end{remark} 

We conclude this subsection with the following central result on Pfaffian systems.
For a smooth map $F\colon \R^n\to\R^k$, a value $y\in\R^k$ is \emph{regular} if $\diff{F}(x)$ is surjective at every $x\in F^{-1}(y)$. A solution of $F(x)=0$ is regular if the differential there has rank $k$; for a square system ($k=n$), we also call such a solution \emph{non-degenerate}.

\begin{theorem}[Khovanskii]\label{thm:khovanskii}
Let $f_1,\dots,f_n\colon\R^n\to\R$ be Pfaffian functions with respect to a common chain of length $s\ge0$ and chain degree $\alpha\ge1$, with component degrees bounded by $\beta_1,\dots,\beta_n\ge1$. Then the number of non-degenerate real solutions of $f_1=\cdots=f_n=0$ is at most
\begin{equation*}
2^{\frac{s(s-1)}2}\beta_1\cdots\beta_n
\left(\beta_1+\cdots+\beta_n-n+\min\{n,s\}\alpha+1\right)^s.
\end{equation*}
\end{theorem}

We use the global version of the theorem, as stated in~\cite[Theorem 3.1]{gabrielov2004complexity}; see also~\cite[\S 3.12, Corollary 5]{khovanskii1991fewnomials}. For general open domains, a bound must also account for the domain's complexity~\cite[Remark 2.12]{gabrielov2004complexity}. Every application below is on all of a Euclidean space. Constant functions can be assigned a positive degree bound when applying this theorem.
A qualitative derivation of finiteness appears in~\cite{Marker1997}, while a self-contained and complete proof appears in~\cite{lotz2026khovanskii}; for parameterwise sharpness of the quantitative bound, see~\cite{bickerton2026sharpnesskhovanskiisbezouttypebound}.

The following component bound follows by the usual compact perturbation and Morse-theoretic argument~\cite[Corollary 3.3]{gabrielov2004complexity}, keeping the $\min\{n,s\}$ term in Theorem~\ref{thm:khovanskii}; see also~\cite{lotz2026khovanskii}.

\begin{corollary}\label{cor:pfaffian-components}
Let $F=(f_1,\dots,f_m)\colon\R^n\to\R^m$ have a common Pfaffian chain of length $s$, chain degree $\alpha\ge1$, and component degrees at most $\beta\ge1$. Then the number of connected components of $\mathcal Z(F)$ is at most
\begin{equation}\label{eq:jones}
2^{\frac{s(s-1)}2+1}\beta(\alpha+2\beta-1)^{n-1}
\left[n(\alpha+2\beta-2)+\alpha(\min\{n,s\}-1)+2\right]^s.
\end{equation}
\end{corollary}

\begin{proof}
As in~\cite[Corollary 3.3]{gabrielov2004complexity}, choose $R$ sufficiently large and then $\eta>0$ sufficiently small so that
\[
F_\eta(x):=\sum_{j=1}^m f_j(x)^2+\eta(\|x\|^2-R^2)
\]
has a smooth compact zero set with at least as many connected components as $\mathcal Z(F)$. A generic linear height function on this hypersurface is Morse and has a critical point on every component. After a rotation, its critical points solve
\[
F_\eta=\partial_2 F_\eta=\cdots=\partial_n F_\eta=0.
\]
All these solutions are non-degenerate. Their degrees are bounded by $2\beta$ and $\alpha+2\beta-1$ for the remaining $n-1$ equations. Substituting these degrees in Theorem~\ref{thm:khovanskii} gives~\eqref{eq:jones}.
\end{proof}
